\section{Definiciones básicas}

\subsection{Algoritmo}

Empezamos suave, con algunas definiciones básicas.

Un \emph{algoritmo} es una secuencia de pasos. Pero nos interesan los algoritmos que resuelven un problema de computación, entonces:

\begin{definicion}{Algoritmo}{algoritmo}
Secuencia de pasos para transformar un conjunto de valores (\emph{entrada}/\emph{\textit{input}}) a un conjunto de valores (\emph{salida}/\emph{\textit{output}}).
\end{definicion}
No se escribe \say{a otro conjunto de valores} porque la salida podría ser igual a la entrada.

Un algoritmo se puede expresar en:
\begin{itemize}
  \item Lenguaje natural (español, inglés, etc.)
  \item Lenguaje formal (matemática o algún lenguaje de programación o lógico)
  \item Seudo código (una mezcla de los dos anteriores)
  \item Como programa de computador
  \item Como hardware (esto es difícil de imaginar en este punto, pero sí)
\end{itemize}
% TODO: Lenguaje formal vs programa de computador, son redundantes?

\subsection{Seudo código}

El seudo código es básicamente algo que se parece a un lenguaje de programación pero que permite:
\begin{itemize}
  \item Despreocuparse de detalles de sintaxis o de implementación, como el manejo de errores.
  \item Usar la forma más clara de expresarse en cada línea: a veces español, a veces matemática, a veces sintaxis de algún lenguaje de programación.
\end{itemize}
Entonces es muy útil para expresar algoritmos.

En los textos formales se establecen convenciones estrictas para el seudo código que se usa; eso elimina la ambigüedad, pero hace que al lector prácticamente le toca aprender un nuevo lenguaje (el seudo código allí definido) para entender el libro. 

En estos apuntes se utiliza un seudo código menos estructurado, basado en la sintaxis de Python, pero con varias libertades, y las ambigüedades las elimino aclarándolas cuando aparecen.

E.g., una función que recibe una lista de numeros y encuentra el mayor:
\begin{minted}{python}
    funcion encontrar_mayor(numeritos: arreglo)
        mayor = numeritos[0]
        for i in 1..numeritos.longitud-1
            if numeritos[i] > mayor
                mayor = numeritos[i]
        return mayor
\end{minted}
Detalles de este seudo código que reflejan que su propósito es ser claro:
\begin{itemize}
  \item Algunas palabras y tipos están en español. A pesar de que no se definieron de antemano, son muy claros porque son los términos que usamos para referirnos a esas cosas.
  \item Omito detalles de sintaxis como los dos puntos, pero preservo los que otorgan claridad, como la indentación.
  \item Evito sintaxis específica de algún lenguaje.
  \begin{itemize}
    \item En el for uso una sintaxis clásica para el rango, $1..n$, que lo hace inambiguo.
    \item Presumo la existencia de un método \texttt{longitud} para obtener la longitud de un arreglo. Es una presunción razonable porque en cualquier lenguaje orientado a objetos existe algo así.
  \end{itemize}
\end{itemize}

\subsection{Problema}

Un problema (en este contexto) es un enunciado que especifica qué se quiere obtener, sin decir cómo. Más formalmente:
\begin{definicion}{Problema (computacional)}{problema}
Especificación de una entrada y de una salida con base en la entrada.
\end{definicion}

E.g., el \emph{problema de ordenamiento} (\textit{sorting problem}, muy famoso):
\begin{itemize}
  \item \textbf{Entrada}: secuencia de $n$ números $a_1, a_2,\ldots,a_n$.
  \item \textbf{Salida}: permutación $a_1', a_2',\ldots,a_n'$ tal que $a_1' < a_2' < \ldots < a_n'$
\end{itemize}

% TODO: Mover esto a un glosario, a una footnote? No es de este curso; es prerrequisito
Una permutación de una secuencia es un posible orden de sus elementos. Si la secuencia tiene más de un elemento, tiene muchas permutaciones. 

Una \emph{instancia} de un problema es una entrada específica, e.g., $1, 4, 2, 3$ es una instancia del problema de ordenamiento.